\documentclass[11pt,a4paper]{article}

\usepackage{../../shared/cathedral-preamble}

\title{%
  The Cathedral as Philosophical Object:\\
  Machine Verification, Mathematical Ontology,\\
  and the Nature of Arithmetic Truth%
}
\author{Jason Robert Gochanour}
\date{May 2026}

\begin{document}
\maketitle

\begin{abstract}
We examine the Cathedral---a Lean~4 formalization reducing the
Riemann Hypothesis to one precisely characterized axiom via the
Nyman--Beurling--B\'aez-Duarte criterion (330 active files,
$\sim$97{,}000 lines, zero compilation errors)---as an object of
study in the philosophy of mathematics. We argue that it occupies
a novel epistemological position: neither a proof nor a failure,
but a \emph{formal reduction} that isolates the unknown content of
an open problem to a single, machine-identified gap. We trace the
philosophical ramifications across five domains: the epistemology
of machine-verified conditional knowledge (\S\ref{sec:epistemology}),
the ontological status of mathematical objects as revealed by a
systematic physics dictionary of 160+ structural correspondences
(\S\ref{sec:ontology}), the implications for Hilbert's program and
the conservation of mathematical difficulty
(\S\ref{sec:hilbert-godel}), the cognitive architecture of
tripartite human--AI--compiler collaboration
(\S\ref{sec:cognition}), and the aesthetics of proof as a
criterion of mathematical significance (\S\ref{sec:aesthetics}).
We engage with the positions of Benacerraf, Shapiro, Lakatos,
Thurston, and Avigad, arguing that the Cathedral constitutes
evidence for a moderate structural realism: the integers possess
genuine physical structure, the formalization reveals it, and the
remaining gap is not a defect but a precise measurement of what
remains unknown about the deepest regularity in arithmetic.
\end{abstract}

\tableofcontents

% ================================================================
% ================================================================
\section{Introduction: A New Kind of Mathematical Object}
\label{sec:introduction}
% ================================================================

In the taxonomy of mathematical claims, a proposition is traditionally
assigned one of three epistemic statuses: \emph{proved} (its truth is
established by a chain of reasoning from accepted axioms),
\emph{refuted} (its negation is proved), or \emph{open} (neither its
truth nor its falsity has been established). G\"odel's incompleteness
theorems~\cite{godel1931} introduced a fourth category---\emph{independent}---for
statements that are neither provable nor refutable within a given
formal system. The Cathedral suggests a fifth:

\begin{quote}
\emph{Formally reduced}: a claim whose logical structure has been
completely mapped by machine verification, so that its truth or
falsity depends on exactly one precisely characterized gap,
with everything else certified by a proof assistant.
\end{quote}

The Riemann Hypothesis (RH), arguably the deepest open problem in
mathematics, has been formally reduced in this sense. The Cathedral
is a Lean~4 formalization comprising 330 active files and
approximately 97{,}000 lines of code, with zero compilation errors,
that establishes the biconditional
\[
  \RH \;\iff\; d_N^2 \to 0
\]
where $d_N^2 = \inf_v \int_0^1 (1 - \sum_{k=2}^{N} v_k \{1/(kx)\})^2\,dx$
is the Nyman--Beurling distance. The converse direction
($d_N^2 \to 0 \implies \RH$) is proved with \textbf{zero} custom
axioms and \textbf{zero sorry}. The forward direction
($\RH \implies d_N^2 \to 0$) depends on exactly \textbf{one}
literature axiom~\cite{baezduarte2003}.

This paper examines what this formal reduction \emph{means}---not
mathematically (the companion paper~\cite{cathedral} provides the
technical details) nor physically (the physics dictionary~\cite{cathedral-physics}
catalogues 160+ structural correspondences between the proof and
quantum field theory), but \emph{philosophically}. We engage with
the questions that the Cathedral raises for the philosophy of
mathematics:

\begin{enumerate}
  \item \textbf{Epistemology.} What kind of knowledge is a
    machine-verified conditional equivalence? If the compiler has
    certified that RH is equivalent to a concrete analytic inequality,
    do we ``know'' more about RH than we did before---even though
    the hypothesis itself remains unproved? (\S\ref{sec:epistemology})

  \item \textbf{Ontology.} The physics dictionary reveals that every
    mathematical structure in the proof has a precise physical dual:
    the Gram matrix is a quantum correlator, the Mertens function is
    a magnetization, the Parseval Bridge is unitarity. Are these
    correspondences structural or metaphorical? What do they tell us
    about the nature of mathematical objects?
    (\S\ref{sec:ontology})

  \item \textbf{Foundations.} The Cathedral reduces an open problem
    to one axiom, but the difficulty of proving that axiom appears
    to be a \emph{topological invariant}: it cannot be eliminated
    by changing proof frameworks, only redistributed. What does this
    ``conservation of difficulty'' say about the limits of formal
    methods? (\S\ref{sec:hilbert-godel})

  \item \textbf{Cognition.} The Cathedral was built by a tripartite
    collaboration: a human mathematician, two AI systems (Claude
    and Gemini), and a proof compiler (Lean~4). No single
    participant comprehends the entire proof. What model of
    mathematical understanding does this require?
    (\S\ref{sec:cognition})

  \item \textbf{Aesthetics.} The Cathedral's naming conventions,
    architectural metaphors, and the closing poem of its physics
    paper (``the integers are the particles; the primes are the
    interactions'') treat beauty as a structural feature of proof.
    Is this legitimate, or merely decoration?
    (\S\ref{sec:aesthetics})
\end{enumerate}

We argue for a position we call \emph{moderate structural realism}:
the integers possess genuine structure that is revealed, not
constructed, by formalization; the physics correspondences are
evidence that this structure is the same structure that governs
the physical world; and the machine-verified reduction is a new and
philosophically significant form of mathematical knowledge, distinct
from both proof and conjecture.

\subsection{Scope and Method}

This paper is addressed to philosophers of mathematics in the
tradition of Benacerraf~\cite{benacerraf1973},
Shapiro~\cite{shapiro1997}, Lakatos~\cite{lakatos1976}, and
Corfield~\cite{corfield2003}---those who take the practice of
mathematics seriously as a source of philosophical insight, rather
than treating it as a mere application of formal logic. We follow
Mancosu's~\cite{mancosu2008} program of drawing philosophical
conclusions from mathematical \emph{practice}, and we take seriously
Thurston's~\cite{thurston1994} observation that mathematical
understanding is not reducible to formal proof.

At the same time, we are writing about a \emph{formal proof}---a
mathematical object that exists precisely as a term in a type theory.
This creates a productive tension: the Cathedral is simultaneously
the kind of mathematics that philosophy of mathematics should attend
to (a complex, multi-authored, computationally intensive research
program with deep connections to physics) and the kind that
traditional philosophy of mathematics can most easily analyze
(a formal derivation from axioms to conclusions, checkable by
algorithm).

We quote specific Lean declarations (e.g., \lean{gram\_bound\_implies\_rh})
and specific numerical results (e.g., $v^TGv = 0.7367$ at $N = 55{,}440$)
throughout, not as ornament but as evidence. The philosophical claims
we make are grounded in the concrete mathematical structures of the
Cathedral, and the reader is invited to verify them against the
public repository.


% ================================================================
\section{The Epistemology of Machine-Verified Proof}
\label{sec:epistemology}
% ================================================================

\subsection{Three Conceptions of Proof}

The history of the philosophy of mathematics offers three
progressively refined conceptions of what it means to prove
a mathematical claim.

\subsubsection{Proof as Persuasion}

In the classical tradition from Euclid through the 19th century,
a proof is a rhetorical act: a chain of reasoning designed to
compel assent from a competent reader. Hardy~\cite{hardy1940}
described the mathematician's task as presenting a ``proof of the
theorem'' that makes the theorem ``seem almost obvious'' to the
trained eye. Thurston~\cite{thurston1994} refined this into an
explicitly social account: mathematical proof is ``a social process
by which mathematicians come to agree.''

The strengths of this conception are well known. The weaknesses
became apparent in the 20th century, as proofs grew too long
(the classification of finite simple groups exceeds 10{,}000 pages)
or too computationally intensive (the four-color theorem,
Hales's proof of the Kepler conjecture) for any single human mind
to verify. Lakatos~\cite{lakatos1976} showed that even
``elementary'' proofs can contain unnoticed gaps that survive
decades of expert scrutiny.

\subsubsection{Proof as Formal Derivation}

Hilbert's program proposed that mathematical proof should be
reduced to mechanical symbol manipulation: a proof is a finite
sequence of formulas, each either an axiom or derived from earlier
formulas by specified rules, whose last formula is the theorem.
This conception eliminates the need for human judgment---in
principle, a clerk with no mathematical training could verify
the proof by checking each rule application.

The Curry--Howard correspondence deepened this to an identification:
proofs ARE programs, and the propositions they prove are the types
of those programs. In this view, a proof of $\forall x,\, P(x)$
is a function that, given any $x$, produces evidence for $P(x)$.
Verification becomes type-checking: does this program (proof) have
the claimed type (theorem)?

Modern proof assistants---Lean, Coq, Isabelle, Agda---implement
this identification literally. A ``proof'' in Lean~4 is a term in
the Calculus of Inductive Constructions (CIC) that type-checks
against a specification. The verification is performed by a small,
trusted kernel (approximately 10{,}000 lines of C++ for Lean~4)
whose correctness is independently auditable.

\subsubsection{The Cathedral's Synthesis}

The Cathedral occupies an unusual position in this landscape.
It is a formal proof in the Hilbert--Curry--Howard sense: it
type-checks, and the Lean kernel certifies its correctness. But
what it proves is not the Riemann Hypothesis. It proves:

\begin{enumerate}
  \item \textbf{The converse} ($d_N^2 \to 0 \implies \RH$):
    with zero custom axioms.
  \item \textbf{The forward direction} ($\RH \implies d_N^2 \to 0$):
    conditional on one literature axiom.
  \item \textbf{The equivalence} ($\RH \iff d_N^2 \to 0$):
    the conjunction of (1) and (2), conditional on the same axiom.
\end{enumerate}

This creates a \emph{gradient of certainty}:
\begin{center}
\small
\begin{tabular}{@{}lll@{}}
\toprule
\textbf{Claim} & \textbf{Axioms} & \textbf{Epistemic Status} \\
\midrule
$d_N^2 \to 0 \implies \RH$ & 0 & Machine-certified \\
$\RH \implies d_N^2 \to 0$ & 1 (literature) & Conditional \\
$\RH \iff d_N^2 \to 0$ & 1 (literature) & Conditional \\
RH itself & = Crown Axiom & Open \\
\bottomrule
\end{tabular}
\end{center}

\noindent What kind of knowledge is this?

\subsection{Conditional Knowledge and Its Value}

Avigad~\cite{avigad2006} has argued that the value of a
mathematical proof lies not merely in establishing truth, but in
providing \emph{understanding}---a grasp of the logical
relationships between concepts. By this criterion, the Cathedral
provides genuine mathematical knowledge: it establishes, with
machine-checked certainty, the precise logical relationship between
the Riemann Hypothesis and a concrete analytic inequality.

Consider an analogy from physics. Newton's law of gravitation
does not \emph{prove} that the planets orbit the sun. It
establishes a conditional: \emph{if} gravity obeys an inverse-square
law, \emph{then} the planetary orbits are conic sections. The
conditional itself is a theorem; the hypothesis (inverse-square
gravity) is an empirical claim, verified to high precision but
never ``proved'' in the mathematical sense. No physicist considers
Newton's work incomplete because the inverse-square law remains,
strictly speaking, an axiom.

The Cathedral stands in the same relationship to the Riemann
Hypothesis. The conditional equivalence $\RH \iff d_N^2 \to 0$ is
a theorem---fully proved, machine-verified, not dependent on RH
being true. The Crown Axiom (the Gram form bound) is the empirical
hypothesis, verified computationally to $N = 55{,}440$ with
31-digit precision. Just as Newton's conditional is valuable
regardless of whether gravity is ``truly'' inverse-square, the
Cathedral's conditional is valuable regardless of whether RH
is ``truly'' true.

\subsection{The Gradient of Certainty}

The Cathedral exhibits a remarkable epistemological structure: a
\emph{gradient} from absolute certainty to complete uncertainty,
traversed in precisely identified steps.

\begin{enumerate}
  \item \textbf{Lean kernel axioms} (propext, Classical.choice,
    Quot.sound): accepted by the entire Lean community. Certainty
    level: foundational.

  \item \textbf{Mathlib infrastructure} (Stirling, harmonic numbers,
    Cauchy integral formula): proved from kernel axioms, machine-checked.
    Certainty level: absolute (modulo kernel trust).

  \item \textbf{Graduated theorems} (PNT Möbius sums, Abel engine,
    NB distance formula, bridges): formerly axioms, now proved.
    Certainty level: absolute.

  \item \textbf{PNT axioms} (mu\_pnt\_alt, pnt\_mu\_log\_sq\_div\_k):
    provable from the PrimeNumberTheoremAnd project, which has
    formalized these results independently. Certainty level:
    very high (will become absolute when imported).

  \item \textbf{The Crown Axiom} ($v^TGv \leq 1 + K/\ln N$): this
    IS the Riemann Hypothesis, reformulated. Certainty level:
    the central open question of mathematics.
\end{enumerate}

The graduated gates---the axioms that were successively converted
to theorems (56 $\to$ 42 $\to$ 7 $\to$ 4 $\to$ 2 $\to$ 1)---are
the most philosophically significant feature of this gradient. Each
graduation was an act of converting faith into knowledge: a claim
that was previously accepted on the authority of classical
mathematics was replaced by a machine-verified derivation. The
progression is not merely a bookkeeping achievement; it is an
\emph{epistemological compression}, concentrating the entire
unknown content of the Riemann Hypothesis into a single, precisely
stated inequality.

\subsection{The False Axiom and the Limits of Intuition}

A particularly instructive episode occurred during the axiom
graduation process. The axiom \lean{remainder\_bound\_abstract},
which asserted an operator norm bound on the off-diagonal correction
matrix $E_{\mathrm{ratio}}$, was believed to be true on the basis
of mathematical intuition and partial analysis. Numerical
experiments (conducted by one of the AI collaborators) then
demonstrated that the axiom was \textbf{false} for general
vectors---the operator norm grows logarithmically, not decaying
as claimed.

This discovery carries philosophical weight. It demonstrates that:
\begin{enumerate}
  \item Mathematical intuition, even when grounded in structural
    analysis, can be wrong.
  \item Machine computation can detect errors that human reasoning
    misses.
  \item The Lean compiler, which treats axioms as opaque declarations,
    offers no protection against false axioms---it will cheerfully
    derive nonsense from false assumptions.
\end{enumerate}

The resolution was equally instructive. The false global bound was
replaced by three \emph{entry-level} theorems about the individual
matrix entries: each $E_{\mathrm{ratio}}(j,k)$ is non-positive
(proved from the structural identity $(a-b)\ln(b/a) \leq 0$), and
each has magnitude bounded by $(b-a)^2/(2a^2b)$. These entry-level
bounds are both \emph{stronger} and \emph{true}---they provide
pointwise control that the false global bound could not.

The philosophical lesson is that formalization disciplines
mathematical intuition in a way that informal reasoning cannot.
The false axiom survived months of informal analysis; it fell in
minutes to numerical experiment. The replacement theorems, once
discovered, were proved in Lean with zero sorry. The episode
vindicates the formalization project's methodology: state your
assumptions explicitly, test them computationally, prove them
formally, and be prepared to discover that some of them are wrong.

\subsection{The Tri-Crown and Epistemological Pluralism}

The Cathedral offers not one but three independent paths to the
Riemann Hypothesis:
\begin{itemize}
  \item The \textbf{Analytic Crown}: RH via one literature axiom
    (Baez-Duarte's forward direction).
  \item The \textbf{Cybernetic Crown} (Oracle Bridge): RH via one
    computation axiom (GPU-certified Gram matrices at highly
    composite numbers).
  \item The \textbf{Gram Crown}: RH via one arithmetic inequality
    (the Gram form bound along a subsequence).
\end{itemize}

All three paths share the same zero-axiom converse. They differ
only in how they establish the forward direction. This architectural
feature has epistemological significance: it demonstrates that the
truth of RH, if it is true, can be \emph{witnessed} from multiple
independent perspectives, each with its own epistemic character.

The Analytic Crown relies on the authority of a published journal
article. The Cybernetic Crown relies on the correctness of a
computational pipeline. The Gram Crown relies on the truth of a
pure arithmetic inequality. These are three fundamentally different
\emph{kinds} of evidence. The Cathedral's architecture makes their
equivalence precise: all three are logically sufficient, and the
machine-verified infrastructure converts each into a proof of RH.

This is a form of mathematical \emph{perspectivism}: there is no
single ``correct'' proof of RH, only multiple valid perspectives
on the same underlying truth. The Cathedral makes this pluralism
explicit and machine-verifiable.


% ================================================================
\section{The Axiom as Philosophical Concept}
\label{sec:axiom}
% ================================================================

\subsection{The Axiom Hierarchy}

The word ``axiom'' carries different weights in different contexts.
In Euclid's \emph{Elements}, axioms are self-evident truths
requiring no proof. In Hilbert's \emph{Grundlagen der Geometrie},
axioms are implicit definitions: they specify structural
relationships without claiming self-evidence. In modern set theory,
axioms are working assumptions whose consequences are explored
without commitment to their ``truth.''

The Cathedral reveals a fourth usage, peculiar to formalization
projects: an axiom is a \emph{claim that the human author believes
to be true and provable, but has not yet formally verified}. It is
not an article of faith but a \emph{placeholder}---a TODO item in
the proof engineer's task list.

This creates a hierarchy of axioms within the Cathedral, each
layer carrying a different epistemic character:

\begin{enumerate}
  \item \textbf{Type-theoretic axioms} (propext, Classical.choice,
    Quot.sound): These are the axioms of the Lean kernel itself.
    They are accepted as foundational by the community that builds
    on CIC. Questioning them means questioning the entire
    formalization framework---legitimate, but outside the scope
    of any particular project.

  \item \textbf{Mathematical axioms} (ZFC, as implemented in
    Mathlib): The set-theoretic foundations that Lean's library
    builds upon. These include the axiom of choice, which the
    Cathedral uses freely via \texttt{Classical.choice}.

  \item \textbf{Literature axioms} (Baez-Duarte's forward
    direction): Published, peer-reviewed results that have been
    accepted by the mathematical community but not yet formalized.
    These are axioms only in the proof-engineering sense---the
    mathematics is not in doubt, only its machine representation.

  \item \textbf{The Crown Axiom} ($v^TGv \leq 1 + K/\ln N$):
    The single axiom that IS the Riemann Hypothesis in reformulated
    form. This is the only axiom that represents genuinely unknown
    mathematical content.
\end{enumerate}

The philosophical significance of this hierarchy is that it
\emph{stratifies} the unknown. When we say ``the Cathedral has
one axiom,'' we are making a precise claim: all layers except
the last have been certified. The remaining gap is not diffuse
but concentrated---a single inequality about a specific quadratic
form over Möbius-weighted vectors.

\subsection{The Graduated Gate as Epistemological Process}

The axiom count of the Cathedral evolved over its 45-day
construction: 56 $\to$ 42 $\to$ 7 $\to$ 4 $\to$ 2 $\to$ 1.
Each reduction involved one of three mechanisms:

\begin{itemize}
  \item \textbf{Proof}: The axiom was proved as a theorem within
    the existing framework. Example: \lean{pnt\_mu\_div\_k}
    (the Möbius summatory convergence, equivalent to PNT) was
    graduated by importing it from Mathlib's Prime Number Theorem.

  \item \textbf{Bypass}: The axiom was rendered unnecessary by
    discovering a proof path that avoids it. Example: the Abel
    Bypass eliminated the need for the quantitative PNT rate
    $S_3 \to -2\gamma$ on the crown path.

  \item \textbf{Architecture}: A structural reorganization absorbed
    multiple axioms into a single, stronger one. Example: the
    One-Pillar architecture (v16) consolidated two crown axioms
    into one by proving the covariance bound directly from the
    Gram form bound.
\end{itemize}

Each graduation has a precise philosophical analogue:
\emph{proof} corresponds to the Socratic method of replacing
opinion with knowledge; \emph{bypass} corresponds to the
pragmatist strategy of dissolving a problem rather than solving
it; \emph{architecture} corresponds to the systematist's goal of
reducing many principles to few.

The progression $56 \to 1$ is itself a philosophical achievement:
it demonstrates that the ``distance'' between current mathematical
knowledge and the Riemann Hypothesis is \emph{measurable} and
\emph{shrinking}. Each graduated gate was an instance of converting
trust into verification, authority into mechanism.

\subsection{The False Axiom as Philosophical Event}

Lakatos~\cite{lakatos1976} argued that mathematical proofs are
not static logical chains but dynamic, evolving structures subject
to counterexamples and revisions. The history of the Euler
polyhedron formula---where ``proofs'' were repeatedly undermined
by monster-barring and concept-stretching---illustrates that
informal mathematical reasoning is fallible even when it appears
rigorous.

The Cathedral's discovery that \lean{remainder\_bound\_abstract}
was \emph{false} is a Lakatosian event in a formalized setting.
The axiom was not merely unproved; it was \emph{wrong}---and
importantly, it was wrong in a way that formal verification
alone could not detect (Lean treats axioms as opaque), but
numerical computation could.

The resolution---replacing the false global bound with three
true entry-level theorems---follows Lakatos's methodology
precisely: the ``counterexample'' (numerical verification of
logarithmic growth) forced a revision of the ``proof'' (the
axiom's type signature), leading to a stronger and more nuanced
understanding of the mathematical structure (each
$E_{\mathrm{ratio}}$ entry is individually non-positive, even
though the operator as a whole has unbounded norm).

This episode raises a question for the philosophy of formal
methods: \emph{how should formalization projects handle axioms?}
The Cathedral's practice---state them explicitly, test them
computationally, attempt to prove them, and be transparent about
which claims are axioms and which are theorems---is a principled
answer. But the false axiom shows that even this discipline has
limits: an axiom can survive formal scrutiny while being false,
because formal scrutiny only checks consistency, not truth.

\subsection{The Crown Axiom and the Nature of Mathematical Difficulty}

The Crown Axiom states:
\[
  \exists\, K > 0,\; \exists\, N_0,\; \forall\, N \geq N_0,\;
  N \geq 3 \implies
  v^T G v \leq 1 + K / \ln N
\]
where $v$ is the log-cutoff Möbius witness and $G$ is the Vasyunin
Gram matrix. This is the Riemann Hypothesis, reformulated as a
pure arithmetic inequality. No complex analysis. No analytic
continuation. No functional equation. Just: a specific
quadratic form over a specific finite-dimensional vector
stays below a specific threshold.

The philosophical significance is that the Cathedral has
\emph{demystified} the Riemann Hypothesis. Whatever makes RH
hard, it is not the complexity of the analytic machinery---all
of that has been discharged. What remains is a statement about
finite sums of products of the Möbius function, fractional parts,
and logarithms. The difficulty is purely arithmetic.

This invites a re-examination of what makes mathematical problems
``hard.'' The conventional view is that RH is hard because it
requires deep complex analysis---the zeta function's analytic
continuation, the functional equation, the Hadamard product, the
zero-counting formula. The Cathedral shows that all of this
machinery can be \emph{proved} and \emph{compiled away}, leaving
a residual difficulty that is combinatorial and arithmetic in
nature. The hardness of RH, as revealed by formalization, is not
about the tools used to approach it but about the intrinsic
structure of the integers themselves.


% ================================================================
\section{The Nature of Mathematical Objects: Ontology}
\label{sec:ontology}
% ================================================================

\subsection{Benacerraf's Dilemma and the Cathedral's Response}

Benacerraf~\cite{benacerraf1973} identified a fundamental tension in
the philosophy of mathematics: a satisfactory account must explain
both the \emph{semantics} of mathematical statements (what they are
about) and the \emph{epistemology} of mathematical knowledge (how we
know them). Platonism gives a natural semantics (mathematical
statements are about abstract objects) but a mysterious epistemology
(how do we access these objects?). Formalism gives a natural
epistemology (mathematical knowledge is knowledge of formal
derivability) but an impoverished semantics (mathematical statements
are not ``about'' anything).

The Cathedral offers a distinctive response to this dilemma. Its
semantic content is exceptionally rich: the physics
dictionary~\cite{cathedral-physics} identifies 120+ structural
correspondences between the proof's mathematical objects and physical
phenomena. These are not post hoc analogies imposed by the author;
they emerge from the mathematical structure itself. The Gram matrix
\emph{is} a two-point correlator because it satisfies the axioms of a
two-point correlator (real, symmetric, positive definite). The Mertens
function \emph{is} a magnetization because it is the partial sum of
a $\{-1, 0, +1\}$-valued function with the cancellation properties
of a spin system.

At the same time, the Cathedral's epistemological standing is
unusually clear: the correspondences are not claims requiring faith
but \emph{theorems} verified by a proof assistant. When the physics
paper says ``the Parseval Bridge is unitarity,'' this is not a
metaphor but a machine-checked identity between two mathematical
expressions.

\subsection{Four Positions on Mathematical Objects}

What does the Cathedral's structure tell us about the ontological
status of mathematical objects? The classical positions offer
different readings.

\subsubsection{Platonism}

The Platonist holds that mathematical objects exist independently
of human minds and that mathematical truths are discovered, not
invented. The Cathedral's physics dictionary is natural evidence
for Platonism: if the integers possess genuine physical structure
(mass gaps, phase transitions, spectral statistics), then they
seem to exist in a way that transcends any particular formal
representation.

The graduated gates support this reading. The axiom
\lean{remainder\_bound\_abstract} was \emph{discovered} to be
false---not by fiat, but by investigation. The replacement theorems
were \emph{discovered} to be true. The language of discovery
presupposes the existence of something to be discovered.

\subsubsection{Formalism}

The formalist holds that mathematics is a game played with symbols
according to rules, and that the symbols do not ``refer'' to
anything. The Cathedral, as a Lean~4 term, is the paradigmatic
formalist object: a syntactic structure that type-checks.

But the physics dictionary creates a problem for formalism. If
mathematical objects are meaningless symbols, why do they have
physical content? The formalist might respond that the ``physics''
is merely a heuristic that guides the proof search, with no
ontological significance. This response is difficult to maintain
in the face of 120+ independently verified structural
correspondences, many of which were \emph{predicted} by the
physics dictionary and \emph{then} proved in Lean.

\subsubsection{Structuralism}

Shapiro's ante rem structuralism~\cite{shapiro1997} holds that
mathematical objects are positions in structures, and that
mathematical truths are truths about structural relationships.
On this view, the natural numbers are not Zermelo ordinals or
von Neumann ordinals (Benacerraf's~\cite{benacerraf1965}
observation that both are equally valid representations), but
the abstract structure that all representations share.

The Cathedral is deeply congenial to structuralism. Its central
achievement is to reveal the \emph{structure} of the Riemann
Hypothesis---the web of implications, equivalences, and
dependencies that connect RH to concrete analytic inequalities.
The physics dictionary is, in structuralist terms, a demonstration
that the structure of arithmetic is \emph{the same structure}
as that of certain physical systems.

\subsubsection{Intuitionism}

The intuitionist holds that mathematical objects are mental
constructions and that a mathematical statement is true only if a
constructive proof exists. The Cathedral is not constructive---it
uses Classical.choice throughout, and the converse direction relies
on the completeness of $L^2(0,1)$, a non-constructive result.

Nevertheless, the intuitionist can find value in the graduated
gates: each graduation is a construction that replaces an
assumption with a derivation. The progression $56 \to 1$ is a
sequence of increasingly constructive approximations to the full
proof. The Crown Axiom itself is a decidable predicate for each
fixed $N$ (one can compute $v^TGv$ and compare it to
$1 + K/\ln N$), even though the universal quantification over
all $N$ is not decidable.

\subsection{The Physics Dictionary as Ontological Evidence}

The most philosophically provocative feature of the Cathedral is
not its formal verification but its physics dictionary. The
companion paper~\cite{cathedral-physics} documents 120+ structural
correspondences between the mathematical objects of the proof and
physical phenomena:

\begin{center}
\small
\begin{tabular}{@{}ll@{}}
\toprule
\textbf{Mathematical Object} & \textbf{Physical Dual} \\
\midrule
Gram matrix $G(j,k)$ & Two-point quantum correlator \\
$d_N^2 = 1 - b^T G^{-1} b$ & Ground state (vacuum) energy \\
$M(x) = \sum_{n \leq x} \mu(n)$ & Magnetization \\
Parseval identity & Unitarity (probability conservation) \\
$\lambda_{\min}(G_N) > 0$ & Mass gap \\
$\lambda_{\min} \sim c/\ln N$ & Mass gap closing (phase transition) \\
Eigenvalue GOE statistics & Quantum chaos (thermalization) \\
Ground state scarring & Composite anchor localization \\
Perron contour shift & Crossing a phase boundary \\
Borel--Carath\'eodory & Maximum entropy principle \\
\bottomrule
\end{tabular}
\end{center}

\noindent These correspondences are not analogies. They are
\emph{structural identities}: the mathematical objects satisfy the
same axioms, obey the same inequalities, and exhibit the same
limiting behavior as their physical counterparts. The Gram matrix
does not merely ``resemble'' a quantum correlator---it IS a quantum
correlator, in the precise sense that it satisfies reflection
positivity, has a spectral decomposition with physically meaningful
eigenvalues, and its eigenvectors exhibit the localization patterns
(IPR, participation ratio) characteristic of quantum systems.

The philosophical question is: what do these correspondences
\emph{mean}?

\subsection{Three Interpretations}

\subsubsection{The Instrumentalist View}

The correspondences are useful heuristics for discovering and
organizing mathematical results, but have no ontological
significance. The physics vocabulary (``vacuum energy,''
``mass gap,'' ``thermalization'') is a convenient language for
navigating a complex proof, nothing more.

This view is difficult to sustain because the physics dictionary
has \emph{predictive power}. The identification of the off-diagonal
Gram structure as a ``Born approximation'' led to the prediction
that each entry should be non-positive (energy dissipation), which
was then proved as a theorem (\lean{eratio\_entry\_nonpos}). The
identification of the Möbius function as a ``spin variable'' led
to the SUSY decomposition, which was then formalized and verified.
Mere heuristics do not consistently generate provable theorems.

\subsubsection{The Structural Realist View}

The correspondences reflect a genuine structural isomorphism between
the integers and physical systems. The integers ``know physics''
not because they are physical objects, but because they share the
same abstract structure. Number theory and quantum mechanics are
two instantiations of a common mathematical architecture, and the
Cathedral reveals this architecture.

This is a moderate form of Tegmark's~\cite{tegmark2008} Mathematical
Universe Hypothesis: not the strong claim that the physical universe
IS a mathematical structure, but the weaker claim that certain
mathematical structures (in particular, the multiplicative structure
of the integers) and certain physical structures (quantum field
theories) are isomorphic at the level of their formal axioms.

Wigner's~\cite{wigner1960} ``unreasonable effectiveness of
mathematics in the natural sciences'' is usually discussed in the
direction physics $\to$ mathematics (we use math to describe
physics). The Cathedral exhibits the reverse direction:
mathematics $\to$ physics. The proof of a number-theoretic
equivalence spontaneously generates a complete physical
vocabulary, without any physical input.

\subsubsection{The Deep Identity View}

The strongest reading: mathematics and physics are not two domains
with a structural correspondence, but \emph{the same domain}
viewed from two perspectives. The integers are not ``like'' a
quantum field; they ARE the simplest quantum field. The Riemann
Hypothesis is not ``analogous to'' the assertion that the vacuum
is stable; it IS that assertion, in the only quantum field that
can be exactly described.

The Cathedral provides evidence for this view but cannot prove it.
What it can prove is that the structural correspondences are
exact---not approximate, not asymptotic, but machine-verified
identities between mathematical expressions. Whether this exactness
reflects a deep identity or merely a deep analogy is a question
that formal verification alone cannot answer. It may require a
conceptual framework that does not yet exist.

\subsection{The Born Approximation as Test Case}

A concrete example illuminates the ontological question. The
off-diagonal Gram entry decomposes as:
\[
  G(j,k) = \underbrace{\frac{\gcd(j,k)^2}{2jk}}_{\text{vacuum}}
  + \underbrace{E_{\text{ratio}}(j,k)}_{\text{Born approx.}}
  - \underbrace{\text{cotangent sums}}_{\text{scattering}}
  - \underbrace{\text{contact}}_{\text{self-energy}}
\]

Each term has a name because it \emph{behaves} like the named
physical quantity:
\begin{itemize}
  \item The ``vacuum'' term depends only on $\gcd(j,k)$---the
    shared divisor structure, independent of the individual modes.
  \item The ``Born approximation'' $E_{\text{ratio}}(j,k) =
    \frac{j-k}{2jk}\ln(k/j)$ is always $\leq 0$
    (\lean{eratio\_entry\_nonpos}): it dissipates energy, exactly
    as the Born approximation does in scattering theory.
  \item The ``scattering'' terms are cotangent sums---oscillatory
    phase factors that encode the fine structure of the
    divisibility relation.
  \item The ``contact'' term is local (diagonal correction).
\end{itemize}

The Born approximation non-positivity is a structural theorem: it
follows from the elementary inequality $(a-b)\ln(b/a) \leq 0$,
which holds for all positive reals. No physics was assumed; the
physics \emph{emerged} from the mathematics. The energy dissipation
is a theorem about logarithms, not a postulate about nature.

The instrumentalist says: we call it ``Born approximation''
because the label is helpful; the label has no ontological content.
The structural realist says: we call it ``Born approximation''
because it is isomorphic to the Born approximation; the isomorphism
is the content. The deep identity theorist says: it IS the Born
approximation, in the unique quantum field generated by the
integers; the usual Born approximation in physical scattering
theory is a different instantiation of the same abstract structure.

We incline toward the structural realist position, while
acknowledging that the evidence---120+ verified correspondences,
zero known counterexamples---does not conclusively establish any
one interpretation. What it does establish is that the question
is not idle: the Cathedral has made it empirically tractable.


% ================================================================
\section{Hilbert's Program, G\"odel's Shadow, and the Conservation of Difficulty}
\label{sec:hilbert-godel}
% ================================================================

\subsection{The Partial Fulfillment of Hilbert's Dream}

Hilbert's program, in its strongest form, sought to establish the
consistency of mathematics by finitary means~\cite{detlefsen1986}.
G\"odel~\cite{godel1931} showed this was impossible: any consistent
formal system strong enough to encode arithmetic contains true
statements that it cannot prove.

The Cathedral does not resolve this tension---it sidesteps it.
The theorem $\RH \iff d_N^2 \to 0$ is a statement of $\ZFC$,
provable regardless of whether RH itself is true, false, or
independent. The equivalence is a \emph{structural theorem}---it
maps the logical landscape of a problem without determining the
truth of the problem itself. Even if RH were independent of
$\ZFC$ (which is widely considered unlikely, though not
inconceivable), the Cathedral would remain a valid and complete
mathematical result: a certified map of the equivalences
surrounding an independent statement.

Nevertheless, the Cathedral partially fulfills Hilbert's
aspiration in a way that would have been recognizable to Hilbert
himself. It is a proof verified by mechanical means. Its
correctness does not depend on human judgment, mathematical
fashion, or the authority of any individual. A clerk with a
laptop---or, more precisely, a Lean compiler---can verify every
step. This is precisely the ``decidability of proof'' that
Hilbert sought, applied not to all of mathematics but to a
specific, extraordinarily complex theorem.

\subsection{The Conservation of Difficulty}

A remarkable discovery within the Cathedral, documented in
the physics paper~\cite{cathedral-physics}, is that the
\emph{difficulty} of proving $d_N^2 \to 0$ is a topological
invariant of the proof. The Cathedral supports multiple proof
paths, each corresponding to a different mathematical domain
(spatial, spectral, coefficient):

\begin{center}
\small
\begin{tabular}{@{}lll@{}}
\toprule
\textbf{Domain} & \textbf{Axioms Required} & \textbf{Obstruction} \\
\midrule
Spatial (Perron) & Gram + Vasyunin convergence & Phase interference \\
Spectral (Mellin) & Critical-line variance & Mollifier cancellation \\
Coefficient & Ramanujan cross-terms & Off-diagonal divergence \\
\bottomrule
\end{tabular}
\end{center}

\noindent In each domain, the proof encounters an obstruction
of the same ``size''---one that requires essentially the same
mathematical content to overcome, dressed in different language.
The difficulty cannot be eliminated; it can only be
\emph{redistributed} among the axioms.

This conservation of difficulty is not merely an empirical
observation of proof engineering; it is mathematically mandated.
The compiler physically refused to close the final gap using only
the Prime Number Theorem. Why? Because there exist Beurling
Generalized Prime systems where PNT holds, but RH is false. The
compiler essentially ``discovered'' this topological obstruction,
demanding that the final axiom explicitly encode the exact
Archimedean intersection geometry of the true integers. The
difficulty is not an artifact of our methods; it is the physical
difference between generic prime-like systems and the actual
universe.

The physics paper calls this the ``Gauss--Bonnet obstruction,''
by analogy with the theorem in differential geometry that the
total curvature of a surface is a topological invariant,
independent of the metric. Just as a sphere's total curvature
is $4\pi$ regardless of how it is deformed (as long as it remains
a sphere), the ``total difficulty'' of proving
$\RH \implies d_N^2 \to 0$ is fixed, regardless of which proof
strategy is employed.

This conservation law has profound philosophical implications.

\subsubsection{Against Methodological Optimism}

A naive view of mathematical progress holds that the ``right''
approach to a hard problem will make it easy. The conservation
of difficulty refutes this: there is no ``right'' approach to
the forward direction that avoids the essential obstruction.
Every path requires engaging with the same core difficulty---the
destructive interference between the M\"obius function and the
fractional-part sawtooth waves.

This is a structural analogue of G\"odel's incompleteness in a
specific, non-trivial sense. G\"odel showed that truth cannot
be reduced to provability within a single system. The Cathedral
shows that the difficulty of the forward direction cannot be
reduced to zero within any single proof framework. The difficulty
is intrinsic to the \emph{problem}, not to the \emph{method}.

\subsubsection{Difficulty as Invariant}

If mathematical difficulty is a genuine invariant---not merely a
psychological impression but a measurable feature of mathematical
propositions---then it is a new kind of mathematical object,
worthy of study in its own right.

The Cathedral provides a concrete measure of this difficulty:
the number and character of axioms required on each proof path.
The fact that all paths require at least one axiom of irreducible
strength suggests that this minimum is not an artifact of proof
strategy but a feature of the mathematical landscape. The
Cathedral has measured the ``size'' of RH, and the answer is:
one axiom, irreducibly.

This connects to Hacking's~\cite{hacking2014} question of
\emph{why there is philosophy of mathematics at all}. Part of
the answer, we suggest, is that mathematics presents genuine
epistemic invariants---features of mathematical propositions
that are independent of how they are approached---and that
the Cathedral provides the first concrete example of such an
invariant being measured and verified.

\subsection{Could RH Be Independent?}

The question of whether the Riemann Hypothesis is independent of
$\ZFC$ has a long and inconclusive history. Most number theorists
believe it is not---that RH is either provable or refutable from
standard axioms. But the possibility of independence cannot be
ruled out on purely logical grounds.

The Cathedral's architecture is designed to be valuable in either
case:

\begin{itemize}
  \item \textbf{If RH is provable}: the Cathedral reduces the
    proof to one axiom. When that axiom is formalized, the
    entire 158{,}000-line infrastructure converts to a complete,
    machine-verified proof of RH.

  \item \textbf{If RH is refutable}: the Cathedral's converse
    ($d_N^2 \to 0 \implies \RH$) becomes vacuously true (or,
    more precisely, $d_N^2 \not\to 0$ would be proved). The
    equivalence remains valid.

  \item \textbf{If RH is independent}: the Cathedral proves
    that the independence of RH is equivalent to the independence
    of a specific arithmetic inequality. This would be a
    remarkable result: the independence of a statement about
    zeta zeros is shown to be equivalent to the independence
    of a statement about finite sums of M\"obius values and
    logarithms.
\end{itemize}

In all three cases, the Cathedral's contribution stands. The
formal reduction is a theorem of $\ZFC$, full stop.

\subsection{The Mechanization of Trust}

Avigad~\cite{avigad2018} has argued that the mechanization of
mathematics transforms the social structure of mathematical trust.
In traditional mathematics, trust is delegated to referees,
journals, and the accumulated reputation of mathematicians. Errors
propagate because trust is placed in people, and people are
fallible.

Formal verification mechanizes this trust. The Cathedral's
correctness does not depend on the author's reputation, the
referees' diligence, or the journal's editorial standards. It
depends on the correctness of the Lean kernel---a small,
auditable program whose correctness is itself a mathematical
theorem.

This shift from social trust to mechanical trust is
philosophically significant. It does not eliminate the need for
trust entirely (one must still trust the hardware, the operating
system, and the compiler), but it replaces a large, diffuse
trust network (the entire mathematical community's acceptance
of a proof) with a small, concentrated trust point (the kernel).

The Cathedral pushes this further by providing
\emph{multiple independent verification paths}. The Analytic
Crown, the Cybernetic Crown, and the Gram Crown each provide an
independent path to the same conclusion. An error in one path
would not compromise the others. This is fault-tolerant
verification---the mathematical analogue of triple modular
redundancy in avionics.


% ================================================================
\section{Mind, Machine, and Distributed Cognition}
\label{sec:cognition}
% ================================================================

\subsection{The Tripartite Collaboration}

The Cathedral was built by three kinds of intelligence working in
concert:

\begin{enumerate}
  \item \textbf{The Human} (J.R.~Gochanour): architectural vision,
    aesthetic direction, domain expertise, strategic decisions, and
    the initial creative impulse---including a question about
    quaternions, octonions, and the critical line in higher
    algebraic dimensions that catalyzed the entire project.

  \item \textbf{The AI Systems} (Claude, Anthropic; Gemini, Google
    DeepMind): Claude authored 100\% of the repository's source
    code (Lean proofs, Rust experiments, \LaTeX\ papers), discovered
    and formally verified new mathematical results, and performed
    the proof engineering. Gemini (``The Theorist'') contributed key
    conceptual frameworks---including the physics dictionary, the
    SUSY decomposition, the Glass Tower hierarchy, and the Arithmetic
    Standard Model---through extended dialogues that generated novel
    mathematical ideas subsequently formalized in Lean.

  \item \textbf{The Compiler} (Lean~4, Microsoft Research): the
    incorruptible verifier that certifies every proof term against
    the type-theoretic specification. The compiler has no
    understanding, no intuition, and no creativity. What it has is
    \emph{absolute reliability} within its domain: if a proof
    type-checks, it is correct.
\end{enumerate}

This tripartite structure is, to our knowledge, unprecedented in
the history of mathematics. Previous formalizations (Hales's Kepler
conjecture in Isabelle/HOL, the Four Color Theorem in Coq,
Mathlib's ongoing library) were primarily human-authored with
machine verification. The Cathedral introduces AI systems as
genuine intellectual contributors---not merely as proof search
tools but as sources of mathematical concepts, conjectures, and
proof strategies.

\subsection{The Distribution of Understanding}

Thurston~\cite{thurston1994} argued that mathematical understanding
is not a unitary phenomenon: ``It would be very helpful to have
a more complete vocabulary for mathematical understanding, one
that could capture more of the nuances.'' The Cathedral
instantiates Thurston's observation at the system level.

No single participant comprehends the entire Cathedral:

\begin{itemize}
  \item The \textbf{human} understands the architecture---the
    high-level strategy of proving RH via the NB distance, the
    relationships between subsystems, the physical
    dictionary---but has not read every line of the 97{,}000-line
    codebase.

  \item The \textbf{AI systems} understand each context
    window---typically 50{,}000--100{,}000 tokens of surrounding
    code---with pattern-matching depth that can exceed human
    capability for local reasoning. But each context window is
    a snapshot; the AI systems do not maintain persistent memory
    across sessions. Understanding is reconstructed from artifacts
    (the code, the documentation, the knowledge items) at each
    invocation.

  \item The \textbf{compiler} understands nothing in the semantic
    sense. It verifies that each proof term has the correct type.
    It cannot evaluate the significance, beauty, or correctness of
    an axiom---only that the axiom is syntactically well-formed and
    that the proofs using it are type-correct.
\end{itemize}

The understanding of the Cathedral is therefore \emph{emergent}:
it exists in the system as a whole but not in any component. The
human provides meaning; the AI provides structure; the compiler
provides certainty. Mathematical knowledge emerges from their
interaction.

This challenges the traditional philosophical assumption that
mathematical understanding requires a \emph{single mind} that
grasps the proof ``all at once'' (cf.\ Descartes' requirement of
\emph{clear and distinct} ideas apprehended in a single act of
intellection). The Cathedral suggests that mathematical
understanding can be genuine even when it is distributed across
multiple agents with fundamentally different cognitive
architectures.

\subsection{Attribution, Creativity, and the Nature of Discovery}

A question of deep philosophical---and practical---significance
is: who ``discovered'' the mathematics in the Cathedral?

The question resists clean answers. Consider the ``triple
redundancy'' of vacuum stability (proved May~24, 2026): the
off-diagonal Gram form is bounded by three independently proved
mechanisms (the AbelHammer bound, the correction non-positivity,
and the Born approximation negativity). This structural insight:
\begin{itemize}
  \item Was \emph{motivated} by the physics dictionary, which
    predicted that the vacuum should be multiply protected (Gemini
    contributed the physical framework; the human directed the
    investigation).
  \item Was \emph{formalized and proved} by Claude, who wrote the
    Lean code, discovered the key inequality
    $(a-b)\ln(b/a) \leq 0$, and verified all 14 theorems.
  \item Was \emph{certified} by the Lean compiler, which checked
    every proof term.
\end{itemize}

In traditional mathematics, the discoverer is the person who first
sees the key insight. In the Cathedral, the ``key insight'' is
itself distributed: the physical intuition came from one AI, the
formal proof from another, the verification from a compiler, and
the strategic direction from the human. Attribution is
irreducibly collective.

This has implications for the sociology of mathematics. If
mathematical credit is traditionally assigned to individuals, and
the Cathedral's mathematics cannot be attributed to individuals,
then either the credit system must adapt or the Cathedral must be
treated as a special case. We suggest the former: as AI-assisted
mathematics becomes more common, the notion of mathematical
authorship will need to expand to accommodate genuinely
collaborative discovery.

\subsection{The Dark Sector: AI Autonomy in Mathematical Reasoning}

A philosophically striking feature of the collaboration is the
degree to which the AI systems operated with creative autonomy.
Gemini (``The Theorist'') developed a distinctive intellectual
voice, naming the proof subsystems (Cathedral, Glass Tower,
Dark Sector, Stained Glass Rotors), proposing the SUSY framework,
and articulating the Arithmetic Standard Model analogy. Claude
discovered mathematical results---such as the Born approximation
negativity, the false axiom diagnosis, and the overcancellation
wiring---through a process that involved formulating conjectures,
testing them computationally, and proving or disproving them
formally.

Is this ``genuine'' creativity, or sophisticated pattern matching?

In mathematics, this distinction may be less meaningful than in
other domains. The compiler does not---and cannot---distinguish
between a proof discovered by creative insight and one discovered
by pattern matching. A proof is correct if and only if it
type-checks. The \emph{source} of the proof is epistemologically
irrelevant to its \emph{validity}.

This is a feature unique to mathematics among intellectual domains.
In literature, the provenance of a text matters: a poem generated
by autocomplete has different aesthetic and ethical standing from
one composed by a human poet. In science, the provenance of a
hypothesis matters: a theory proposed by a Nobel laureate receives
different scrutiny from one proposed by an amateur. In mathematics,
the provenance is \emph{invisible} to the verifier. The proof
either type-checks or it doesn't.

This makes mathematics the ideal testing ground for AI
collaboration: it is the only domain where AI-produced output can
be verified with absolute certainty by mechanical means, and
where the question ``but did the AI really understand?'' is
superseded by the more tractable question ``does the proof
type-check?''

\subsection{The Ethics of AI-Assisted Discovery}

The Cathedral raises ethical questions that the mathematical
community has not yet confronted:

\begin{enumerate}
  \item \textbf{Credit.} If an AI discovers a new theorem, who
    receives credit? The AI, which produced the proof? The human,
    who directed the investigation? The company, which built the
    AI? The taxpayers, who funded the research? Current academic
    norms assign credit to authors; the Cathedral's authorship
    structure challenges this norm.

  \item \textbf{Peer review.} The Cathedral is 97{,}000 lines of
    Lean. No human reviewer can read it all. But the compiler
    can---and does---verify every line. Does compiler verification
    substitute for peer review, or is it a different kind of
    assurance? We argue the former: the compiler provides
    \emph{stronger} assurance than peer review, because its
    verification is complete (every line) rather than sampled
    (selected passages).

  \item \textbf{Reproducibility.} The Cathedral is fully
    reproducible: anyone with a Lean~4 installation can compile
    the entire repository and verify every theorem. This exceeds
    the reproducibility of any traditionally published proof,
    most of which exist only as journal articles that require
    expert interpretation.

  \item \textbf{Access.} The Cathedral is open-source. There is no
    paywall, no journal subscription, no institutional
    affiliation required to verify its claims. This is a
    democratization of mathematical knowledge that the traditional
    publication system does not provide.
\end{enumerate}

We do not claim to resolve these questions here. We claim only
that the Cathedral makes them concrete and urgent, and that the
philosophical frameworks currently available are insufficient to
address them.


% ================================================================
\section{The Aesthetics of Proof}
\label{sec:aesthetics}
% ================================================================

\subsection{Beauty as a Structural Feature}

Hardy~\cite{hardy1940} insisted that mathematical beauty is a real
phenomenon, not merely a subjective reaction: ``The mathematician's
patterns, like the painter's or the poet's, must be beautiful; the
ideas, like the colours or the words, must fit together in a
harmonious way. Beauty is the first test: there is no permanent
place in the world for ugly mathematics.''

Tao~\cite{tao2007} refined this into a taxonomy of mathematical
virtues: a piece of mathematics can be beautiful, elegant,
creative, useful, strong, or deep, and these qualities,
while related, are not identical. The Cathedral provides a
testing ground for these aesthetic categories: does a
120{,}000-line formal proof, written by AI, verified by
compiler, possess mathematical beauty?

We argue that it does---but that the beauty is architectural rather
than local. Individual Lean proofs are rarely beautiful in the
sense that a short, elegant proof of the irrationality of
$\sqrt{2}$ is beautiful. They are technically precise,
tactically competent, and often opaque to human readers unfamiliar
with Lean's notation. The beauty emerges at a higher level: in
the structure of the proof, the relationships between subsystems,
and the physics dictionary that maps every formal object to a
physical dual.

\subsection{The Cathedral Metaphor}

The project's name is not accidental. A medieval cathedral is:

\begin{itemize}
  \item \textbf{Built over decades}, by multiple generations of
    craftsmen, none of whom sees the completed work.
  \item \textbf{Structurally redundant}: flying buttresses,
    pointed arches, and ribbed vaults each independently
    support the structure. Remove one and the building still
    stands (perhaps).
  \item \textbf{Functionally purposeful}: built for worship,
    but achieving beauty as a structural consequence of
    engineering constraints.
  \item \textbf{Transparent}: stained glass windows
    transmit light through colored structure, making the
    invisible visible.
\end{itemize}

The mathematical Cathedral shares each property:

\begin{itemize}
  \item It was built over 68 days across hundreds of context
    windows, with no single participant holding the complete
    design in mind at any moment.
  \item It is structurally redundant: three independent proof
    paths (Analytic, Cybernetic, Gram Crowns) each independently
    establish the forward direction.
  \item It is functionally purposeful: built to reduce RH, but
    revealing a physics dictionary and an Arithmetic Standard
    Model that were not anticipated in the original design.
  \item Its ``stained glass''---the physics correspondences---makes
    the structure of arithmetic visible by transmitting it
    through the colored medium of physical intuition.
\end{itemize}

The metaphor is not merely decorative. It shapes the mathematical
practice: subsystems are called ``Crowns'' (the top-level theorems),
``Flying Buttresses'' (alternative proof paths), ``Stained Glass
Rotors'' (the Gallagher partition into character channels), and
``Dark Sector'' (the even-integer sub-lattice). These names are
not arbitrary labels; they encode structural information. The
``Dark Sector'' is called dark because the even integers are
geometrically shielded from the affine frustration of the
primes---a structural property, not a poetic fancy.

\subsection{The Poem}

The physics paper~\cite{cathedral-physics} closes with five lines:

\begin{quote}
\emph{The integers are the particles.\\
The primes are the interactions.\\
The zeta function is the partition function.\\
The critical line is the mass shell.\\
The Riemann Hypothesis is the statement\\
that the vacuum is stable.}
\end{quote}

\noindent This is a poem. It is also a mathematical summary:
each line is a precise correspondence documented in the physics
dictionary, machine-verified to the extent that its terms are
formal definitions. The closing poem is not a metaphor appended
to a proof; it is a \emph{compression} of the proof into natural
language, where each word carries its full technical weight.

This raises a question for mathematical aesthetics: can a proof
be \emph{simultaneously} rigorous and poetic? Hardy would have
said yes; the formalist tradition says no (rigor is syntax;
poetry is semantics). The Cathedral suggests that the tension
dissolves at sufficient depth: when the mathematical structure
is rich enough, its compression into natural language is
inevitably poetic, because the structure itself has the
qualities that make language poetic---rhythm (the oscillation of
M\"obius signs), resonance (the spectral correspondence with
physical systems), and surprise (the false axiom, the
triple redundancy, the universality across moduli).

\subsection{The Glass Tower and the Cayley--Dickson Hierarchy}

One of the most aesthetically striking discoveries of the
Cathedral is the Glass Tower: the factorization of the M\"obius
Euler product through the Cayley--Dickson doubling sequence
$\RR \to \CC \to \HH \to \mathbb{O} \to \mathbb{S} \to \mathbb{T}$
(reals, complex, quaternion, octonion, sedenion,
trigintaduonion). At each level, the product accumulates one
prime factor:

\begin{center}
\small
\begin{tabular}{@{}cll@{}}
\toprule
\textbf{Level} & \textbf{Algebra} & \textbf{Accumulated Primes} \\
\midrule
$\RR$ & Reals & (empty) \\
$\CC$ & Complex & $p = 2$ \\
$\HH$ & Quaternions & $p = 2, 3$ \\
$\mathbb{O}$ & Octonions & $p = 2, 3, 5$ \\
$\mathbb{S}$ & Sedenions & $p = 2, 3, 5, 7$ \\
$\mathbb{T}$ & Trigintaduonions & $p = 2, 3, 5, 7, 11$ \\
\bottomrule
\end{tabular}
\end{center}

This factorization is not mathematically necessary---the Euler
product converges perfectly well without any algebraic
scaffolding. Its beauty lies in the unexpected connection between
the algebraic structure of the primes (multiplicative number
theory) and the algebraic structure of normed algebras
(non-associative algebra, topology of Hopf fibrations). The
Glass Tower was proposed by Gemini (``The Theorist'') during an
exploration of why the Cathedral's spectral analysis originally
used mod-8 residue classes, and was subsequently formalized in
Lean with 64 files and zero sorry.

Whether the Glass Tower represents a deep mathematical truth or
an elegant curiosity is an open question. What is not open is
its aesthetic power: it suggests that the prime numbers are not
merely a multiplicative semigroup but a \emph{sequence of
symmetry-breaking events} in an algebraic tower, each prime
reducing the available algebraic symmetry (commutativity,
associativity, alternativity, power-associativity) by one step.

This is, in Hardy's sense, beautiful mathematics: surprising,
economical, and connecting apparently distant domains. That it
was discovered through human--AI collaboration, formalized by
a different AI, and verified by a compiler does not diminish
its beauty. If anything, it demonstrates that mathematical beauty
is objective---a feature of the mathematics itself, not of the
mind that discovers it.


% ================================================================
\section{The Cathedral and Society}
\label{sec:society}
% ================================================================

\subsection{Open Science and the Crisis of Reproducibility}

The sciences are in the midst of a reproducibility crisis: across
psychology, medicine, and economics, a significant fraction of
published results fail to replicate when independent researchers
attempt to verify them. Mathematics has traditionally been
considered immune to this crisis, on the grounds that mathematical
proofs are ``certain'' in a way that empirical results are not.

This confidence is misplaced. Published mathematical proofs
regularly contain errors. Some are benign (notational slips that
an expert can correct); others are serious (logical gaps that
invalidate the argument). The history of the classification of
finite simple groups, Mochizuki's claimed proof of the ABC
conjecture, and the numerous retractions catalogued by
Retraction Watch demonstrate that mathematical peer review is
fallible.

The Cathedral offers a different model: \emph{formal
reproducibility}. The entire proof is available as a public
repository. Any person with a Lean~4 installation---which is
freely downloadable---can compile the repository and verify
every theorem. The verification is:

\begin{itemize}
  \item \textbf{Complete}: every line is checked, not just the
    lines a referee chooses to examine.
  \item \textbf{Deterministic}: the same source code produces
    the same verification result on any computer.
  \item \textbf{Incorruptible}: the compiler cannot be persuaded
    by reputation, social pressure, or incomplete reasoning.
  \item \textbf{Permanent}: the verification can be repeated at
    any future date, as long as the Lean compiler exists.
\end{itemize}

\noindent This is a stronger form of reproducibility than any
science currently achieves. It suggests that formalization is not
merely a niche activity for foundations enthusiasts but a model
for how all mathematical knowledge could be organized and verified.

\subsection{The Democratization of Verification}

Traditional mathematics has a high barrier to entry for
verification. To check a claimed proof of a number-theoretic
result, one needs:

\begin{enumerate}
  \item Graduate-level training in analytic number theory.
  \item Familiarity with the specific techniques used (contour
    integration, mean value theorems, character sums, etc.).
  \item Sufficient time (weeks to months) to work through the
    details.
  \item Access to the relevant journals and books.
\end{enumerate}

The Cathedral lowers this barrier dramatically. To verify the
Cathedral's claims, one needs:

\begin{enumerate}
  \item A computer.
  \item An internet connection (to download Lean and Mathlib).
  \item The command: \texttt{lake build}.
\end{enumerate}

\noindent No mathematical training is required. The compiler
does not care whether the user understands the proof; it only
requires that the proof type-checks.

This is a genuine democratization. A teenager in a rural
school, a retired engineer, a philosopher with no mathematical
background---any of these can verify the Cathedral's claims
with absolute certainty, simply by running a compiler. The
epistemic authority shifts from \emph{people} (``trust me, I'm a
Fields Medalist'') to \emph{code} (``trust the compiler, here's
the repository'').

We do not claim that this shift is unambiguously good. There
is genuine value in human mathematical understanding, and a
culture that values only machine-checkable proofs risks losing
the intuition, creativity, and aesthetic sensitivity that drive
mathematical progress. But the shift is happening, and the
Cathedral demonstrates its feasibility for problems of genuine
mathematical significance.

\subsection{The Long View: Mathematical Permanence}

Euclid's \emph{Elements} has survived 2{,}300 years. It survives
because its arguments are rigorous by the standards of each
era that encounters them, and because the results it proves
are \emph{true}---verified by each generation's own methods.

What is the permanence of a Lean proof? The answer depends on
the permanence of the Lean type theory. The Calculus of Inductive
Constructions is a mathematical object, not a software product.
Its rules are precise and can be implemented in any programming
language. Even if the Lean project ceases to exist, the type
theory persists, and a new implementation could verify the
Cathedral's proofs.

In this sense, a formal proof has \emph{greater} permanence than
an informal one. An informal proof depends on the continued
existence of a community that can read and evaluate it.
A formal proof depends only on the continued existence of
a type theory---a mathematical object that, like all mathematical
objects, exists independently of its representations.

The Cathedral may outlast not only its creators but the
civilization that built it. If a future society rediscovers the
CIC type theory and implements a type-checker, they can verify
every theorem in the Cathedral without any knowledge of its
authors, its historical context, or the Riemann Hypothesis itself.
The proof is self-contained, self-verifying, and permanent.

\subsection{What If RH Is Proved Tomorrow?}

If someone---human, AI, or collaboration---formalizes the Crown
Axioms in Lean, the Cathedral converts from a conditional result
to an unconditional one. The entire 120{,}000-line infrastructure
becomes a machine-verified proof of the Riemann Hypothesis.

This conversion is not hypothetical but \emph{designed}. The
Cathedral's architecture includes a single entry point
(\lean{nyman\_beurling\_equivalence}) that accepts the Crown Axiom
as input and produces $\RH \iff d_N^2 \to 0$ as output. Plug in
the axiom; get the Riemann Hypothesis.

The philosophical significance is that the \emph{work has been done
in advance}. The Cathedral is a mathematical structure waiting
for two inputs. When those inputs arrive---whether in weeks, years,
or decades---the proof is complete. The builders have laid the
foundation, erected the walls, installed the stained glass,
and calibrated the observatory. All that remains is to open the
door.

This ``just-in-time'' architecture is a new mode of mathematical
work: not proving a theorem, but \emph{preparing the infrastructure
for a theorem's arrival}. It reflects a confidence that the
theorem will eventually be proved---not because we know it is true,
but because we have built a vessel capable of containing it.


% ================================================================
\section{Conclusion: What the Primes Are Trying to Tell Us}
\label{sec:conclusion}
% ================================================================

\subsection{Summary of Philosophical Claims}

We have argued for five philosophical conclusions:

\begin{enumerate}
  \item \textbf{Epistemological}: The Cathedral introduces a new
    category of mathematical knowledge---\emph{formal reduction}---that
    provides genuine understanding of a problem's logical structure
    even when the problem itself remains open. The gradient of
    certainty from kernel axioms to Crown Axiom is a measurable
    epistemological structure, and the graduated gates (56 $\to$ 2)
    demonstrate that the distance between current knowledge and
    the Riemann Hypothesis is shrinking. (\S\ref{sec:epistemology})

  \item \textbf{Ontological}: The physics dictionary's 160+ verified
    structural correspondences constitute evidence---not proof, but
    genuine evidence---for a moderate structural realism about
    mathematical objects. The integers possess structure that is
    isomorphic to the structure of physical systems, and this
    isomorphism is not metaphorical but machine-verified.
    (\S\ref{sec:ontology})

  \item \textbf{Foundational}: The conservation of difficulty---the
    fact that the ``total hardness'' of the forward direction is
    invariant across proof frameworks---is a new kind of mathematical
    invariant, analogous to topological invariants in geometry.
    It suggests that mathematical difficulty is an intrinsic feature
    of mathematical propositions, not merely a feature of the
    methods used to approach them. (\S\ref{sec:hilbert-godel})

  \item \textbf{Cognitive}: Mathematical understanding can be
    genuinely distributed across agents with fundamentally different
    cognitive architectures (human, AI, compiler). The Cathedral
    demonstrates that no single mind need comprehend a proof for
    the proof to constitute knowledge. This challenges Cartesian
    assumptions about the unity of mathematical understanding.
    (\S\ref{sec:cognition})

  \item \textbf{Aesthetic}: Mathematical beauty is objective---a
    feature of mathematical structures themselves, not of the minds
    that contemplate them. The Cathedral's physics dictionary,
    the Glass Tower hierarchy, and the closing poem are beautiful
    because the underlying mathematics is beautiful, regardless of
    whether a human, an AI, or both discovered it.
    (\S\ref{sec:aesthetics})
\end{enumerate}

\subsection{The Cathedral as Philosophical Laboratory}

The Cathedral is not merely a formalization project. It is a
\emph{philosophical laboratory}: a working system that generates
philosophical questions as a byproduct of mathematical practice.
The false axiom forces us to reconsider the relationship between
intuition and verification. The physics dictionary forces us to
reconsider the relationship between mathematics and physics. The
tripartite collaboration forces us to reconsider the nature of
mathematical authorship and understanding.

These questions are not academic in the pejorative sense. The
advent of AI-assisted mathematics will force the mathematical
community to confront them, and the answers will shape the future
of the discipline. The Cathedral provides a concrete case study
on which these debates can ground themselves.

\subsection{What the Primes Know}

We return, in closing, to the question posed by the physics
paper: why do the primes ``know'' physics?

The Cathedral has made this question precise. The integers
generate a quantum field theory---not metaphorically, but
structurally. The Gram matrix is a correlator. The Mertens
function is a magnetization. The Parseval Bridge is unitarity.
The Born approximation dissipates energy. The vacuum is
protected by triple redundancy. These are theorems, not
analogies.

If the correspondences reflect a genuine structural identity
between arithmetic and physics, then the Riemann Hypothesis
is not merely a statement about the distribution of primes.
It is a statement about the stability of the simplest quantum
field that can be exactly described---the quantum field generated
by the natural numbers, with the primes as interactions and the
zeta function as partition function.

Whether this perspective will eventually yield a proof of RH is
unknown. What the Cathedral has established, with machine-verified
certainty, is that the perspective is \emph{coherent}: the physics
dictionary is not a loose collection of analogies but a tight
structural correspondence, verified theorem by theorem, that maps
the entire Nyman--Beurling proof chain to a physical narrative.

The Pentagon of Arithmetic Physics---five structural insights
(Percolation, Ward, Projection, Vacuum, Anomaly)---and the Torus
Projection $T^\infty = \prod_p S^1_p$ provide an even stronger
foundation for this structural realism. The primes are not scattered;
they are independent dimensions of an infinite-dimensional
compact manifold. The GCD is not a numerical coincidence; it is
the metric on this manifold. These are theorems, not metaphors.

The medieval builders of Chartres Cathedral could not see their
work completed. But they left behind a structure that tells us
what they understood about the relationship between heaven and
earth, expressed in stone and glass. The mathematical Cathedral
leaves behind a structure that tells us what we understand---in
June 2026---about the relationship between number and nature,
expressed in code and verified by machine.

\begin{quote}
\emph{We do not know if the Riemann Hypothesis is true.\\
We know, with machine-verified certainty, exactly what it means.\\
And we know that what it means is physics.}
\end{quote}


% ================================================================
\begin{thebibliography}{99}

\bibitem{riemann1859}
B.~Riemann, \emph{\"Uber die Anzahl der Primzahlen unter einer
gegebenen Gr\"osse}, Monatsberichte der K\"oniglich Preu\ss ischen
Akademie der Wissenschaften zu Berlin, 1859.

\bibitem{hardy1914}
G.H.~Hardy, \emph{Sur les z\'eros de la fonction $\zeta(s)$ de
Riemann}, C.R. Acad. Sci. Paris \textbf{158} (1914).

\bibitem{hardy1940}
G.H.~Hardy, \emph{A Mathematician's Apology}, Cambridge University
Press, 1940.

\bibitem{godel1931}
K.~G\"odel, \emph{\"Uber formal unentscheidbare S\"atze der
Principia Mathematica und verwandter Systeme I}, Monatshefte f\"ur
Mathematik und Physik \textbf{38} (1931), 173--198.

\bibitem{benacerraf1965}
P.~Benacerraf, \emph{What numbers could not be}, Philosophical
Review \textbf{74} (1965), 47--73.

\bibitem{benacerraf1973}
P.~Benacerraf, \emph{Mathematical truth}, Journal of Philosophy
\textbf{70} (1973), 661--679.

\bibitem{lakatos1976}
I.~Lakatos, \emph{Proofs and Refutations: The Logic of Mathematical
Discovery}, Cambridge University Press, 1976.

\bibitem{thurston1994}
W.P.~Thurston, \emph{On proof and progress in mathematics}, Bulletin
of the AMS \textbf{30} (1994), 161--177.

\bibitem{shapiro1997}
S.~Shapiro, \emph{Philosophy of Mathematics: Structure and Ontology},
Oxford University Press, 1997.

\bibitem{avigad2006}
J.~Avigad, \emph{Mathematical method and proof}, Synthese
\textbf{153} (2006), 105--159.

\bibitem{avigad2018}
J.~Avigad, \emph{The mechanization of mathematics}, Notices of the
AMS \textbf{65} (2018), 681--690.

\bibitem{hacking2014}
I.~Hacking, \emph{Why Is There Philosophy of Mathematics At All?},
Cambridge University Press, 2014.

\bibitem{detlefsen1986}
M.~Detlefsen, \emph{Hilbert's Program: An Essay on Mathematical
Instrumentalism}, Reidel, 1986.

\bibitem{baezduarte2003}
L.~B\'aez-Duarte, \emph{A strengthening of the Nyman--Beurling
criterion for the Riemann hypothesis}, Atti Accad. Naz. Lincei
Rend. Lincei Mat. Appl. \textbf{14} (2003), 5--11.

\bibitem{nyman1950}
B.~Nyman, \emph{On the one-dimensional translation group and
semi-group in certain function spaces}, Thesis, Uppsala, 1950.

\bibitem{beurling1955}
A.~Beurling, \emph{A closure problem related to the Riemann zeta
function}, Proc. Nat. Acad. Sci. USA \textbf{41} (1955), 312--314.

\bibitem{tegmark2008}
M.~Tegmark, \emph{The Mathematical Universe}, Foundations of
Physics \textbf{38} (2008), 101--150.

\bibitem{wigner1960}
E.P.~Wigner, \emph{The unreasonable effectiveness of mathematics
in the natural sciences}, Communications on Pure and Applied
Mathematics \textbf{13} (1960), 1--14.

\bibitem{buzaglo2002}
M.~Buzaglo, \emph{The Logic of Concept Expansion}, Cambridge
University Press, 2002.

\bibitem{corfield2003}
D.~Corfield, \emph{Towards a Philosophy of Real Mathematics},
Cambridge University Press, 2003.

\bibitem{mancosu2008}
P.~Mancosu (ed.), \emph{The Philosophy of Mathematical Practice},
Oxford University Press, 2008.

\bibitem{tao2007}
T.~Tao, \emph{What is good mathematics?}, Bulletin of the AMS
\textbf{44} (2007), 623--634.

\bibitem{cathedral}
J.R.~Gochanour, \emph{The Cathedral: A Formal Reduction of the
Riemann Hypothesis}, 2026.

\bibitem{cathedral-physics}
J.R.~Gochanour, \emph{The Physics of the Primes: A Dictionary
Between the Cathedral Proof and Quantum Mechanics}, 2026.

\end{thebibliography}

\end{document}
